\subsection{The Role of Empathy and Compassion in Human Social Behavior}
\label{sec:2.2.1}
“Empathy” is asked to cover three distinct things, and the distinction is not academic. Cognitive empathy is understanding another's emotional state — inferring that someone is grieving or anxious, the mentalizing work of theory of mind pointed at feelings rather than beliefs. Affective empathy is actually sharing the state: feeling a version of another's distress, the wince at someone else's slammed door. And compassion — sometimes called empathic concern — is caring about another's welfare and being moved to help, which is a motivation, not a perception or a shared feeling.

Compassion is not, as it is often described, an evolutionary extension of empathy. The psychologists Tania Singer and Olga Klimecki trained one group of people, over weeks, first in empathy practice (attending to and resonating with others' suffering) and then in compassion practice (cultivating warmth and concern in the tradition of loving-kindness meditation), tracking what each stage did to brain and mood against a separate control group given memory training instead \autocite{klimecki2013differential}. Empathy training increased negative affect and strengthened the shared-distress network — the anterior insula and anterior cingulate cortex, the same regions active when a person feels their own pain. Sharing more suffering made people feel worse, exactly as a shared-distress mechanism predicts. Compassion training produced the opposite emotional signature: negative affect did not rise, positive affect did, even while contemplating others' distress, and the brain activity moved to a different network entirely — the medial orbitofrontal cortex and ventral striatum, the circuitry of reward, affiliation, and approach rather than of shared pain \autocite{klimecki2013differential}. The two could be trained separately. Compassion is not more empathy. It is a distinct, trainable system that can take shared distress as an input without being flooded by it — the difference between a clinician wrecked by every patient's suffering and one whose steady warmth lets them act all day without collapse.

A teacher who feels a struggling student's frustration as their own distress is at risk of withdrawing from exactly the student who needs them most.

The mirror neuron system is often invoked as the mechanism by which one person simulates another's experience, and that claim has not survived contact with the evidence \autocite{hickok2014myth}. The episode is worth carrying into machine design. Mirror neurons were found in the early 1990s in macaque premotor cortex, cells that fire both when the animal grasps an object and when it watches someone else do it \autocite{gallese1996action}, and the interpretation outran the finding: they were cast as the neural basis of empathy, language, imitation and cultural transmission at once \autocite{ramachandran2009neurons}. Single-neuron evidence was never obtained in humans on the same terms and the recordings that do exist are markedly less clean \autocite{mukamel2010singleneuron}; the broken-mirror account of autism repeatedly failed to find the reduced activity it predicted \autocite{southgate2008unbroken}. What survives is narrower — some analog probably exists in humans, and it plausibly contributes to recognizing actions already within one's own motor repertoire. Knowing what someone is doing and knowing what they believe are different computations, and the brain appears to run them in different places. A single vivid mechanism, well named, absorbed four distinct capacities that later had to be separated again, and architectures for artificial empathy are exposed to the same failure.

Cognitive and affective empathy both depend on perspective-taking and on emotion regulation — managing one's own response well enough to stay useful. Compassion depends on those too and adds what neither supplies: a motivation to act that does not require manufacturing the other person's distress in oneself first.

For an AI system, this argues against a single “empathy module” that a designer either builds or doesn't. Cognitive empathy is a tractable perception and inference problem. Affective empathy, engineered as an internal state that mirrors the user's, buys the system nothing on its own and imports the burnout risk that afflicts overtaxed human carers. What a helping system actually needs is the third thing: a stable disposition to act on another's welfare that does not depend on the system first constructing a version of that person's suffering in itself.
